\par\Needspace{7\baselineskip}\bigskip
\begin{center}
\pdfbookmark[2]{Gauss an Olbers, Göttingen, 15. Juni 1814}{pallas-brief-18140615}
\textsc{Gauss} an \textsc{Olbers}, Göttingen, 15. Juni 1814.
\end{center}
\noindent\hspace*{\parindent}. . . . .\quad Mit der Pallas geht es wunderlich. \textsc{Nicolai} hat die numerische Berechnung der Saturnsstörungen für die 9 Oppositionen vollendet; aber die vorherige schöne Übereinstimmung wird dadurch etwas schlechter. Immer aber bleibt die Nothwendigkeit der Vergrösserung der Jupitersmasse entschieden, und wird noch etwas grösser, nemlich \(=\frac{1}{40}\mathpunct{\raisebox{.3ex}{.}}\) Ohne Verbesserung der Jupitersmasse wird die möglich beste Übereinstimmung der Längen wie folgt:
\begin{center}
\begingroup\small\setlength{\arraycolsep}{3pt}
\(\underbrace{\begin{array}{@{}rrrrrrrrr@{}}
-6''{,}7&-27''{,}7&-4''{,}2&+28''{,}7&+31''{,}0&-18''{,}7&-6''{,}5&-33''{,}3&+0''{,}1
\end{array}}\)\par
Summe der Quadrate . . . . . . 4110\par\smallskip
Mittlerer Fehler \(=\BareRoot{\frac{4110}{9}}=21''\).
\endgroup
\end{center}
\par Mit Vergrösserung der Jupitersmasse um \(\frac{1}{40}\) hingegen:
\begin{center}
\begingroup\small\setlength{\arraycolsep}{3pt}
\(\underbrace{\begin{array}{@{}rrrrrrrrr@{}}
-2''{,}3&-9''{,}8&+11''{,}3&+8''{,}5&+5''{,}1&-8''{,}8&-5''{,}7&-5''{,}7&+7''{,}4
\end{array}}\)\par
Summe der Quadrate \(=525\), die sich aber durch eine kleine Nachfeilung noch auf 503 herunterbringen lässt, so dass mittlerer Fehler \(=7''{,}5\) wird.
\endgroup
\end{center}
\par Die Breiten hingegen stimmen sehr schön und noch etwas besser, als vor Zuziehung der {\AstroSymbols ♄}-Störungen, nemlich:
\begin{center}
\begingroup\small\setlength{\tabcolsep}{3pt}
\begin{tabular}{@{}rrrrrrrrr@{}}
\(+5''{,}8\)&\(-0''{,}1\)&\(-2''{,}0\)&\(-3''{,}2\)&\(-8''{,}5\)\footnotemark&\(-1''{,}7\)&\(-8''{,}2\)&\(+1''{,}2\)&\(-4''{,}3\).
\end{tabular}
\endgroup
\end{center}
\footnotetext{Die schlecht beobachtete {\AstroSymbols ☍} in der Breite.}
\par\SourceLeaf{431}Ob bei den Längen vielleicht ein Fehler bei einem Elemente in Hrn. \textsc{Nicolais} Rechnung eingeschlichen ist, will ich nicht entscheiden; sonderbar ist es freilich, dass, wenn ich allen Störungen der Epoche das entgegengesetzte Zeichen gebe, die Übereinstimmung \textls[180]{weit} besser wird, und die Summe der Quadrate der Fehler nur etwa 200 beträgt; ich habe Hrn. \textsc{Nicolai} eine Methode angegeben, dies zu prüfen. Die Marsstörungen werde ich schwerlich früher als binnen \(\frac{1}{2}\) Jahr endigen, ob sie gleich sehr klein sind; mein Eigensinn ist einmal, sie so zu absolviren, dass nichts zu wünschen übrig bleibt.
\par
